\SourceLeaf{042}
\[
\begin{aligned}
K+\kappa&=\theta(fF+gG+hH+\text{etc.})+K,\quad\text{sive}\\[3pt]
\theta&=\frac{\kappa}{\omega}
\end{aligned}
\]
accipiendo \(\omega\) in eadem significatione ut in art. praec. Et quum \(\Omega-M\), indefinite, sit functio homogenea secundi ordinis indeterminatarum \(\xi,\eta,\zeta\) etc., sponte patet, eius valorem pro \(\xi=\theta f,\allowbreak\eta=\theta g,\allowbreak\zeta=\theta h\) etc. fieri \(=\theta\theta\omega\), et proin valorem minimum, quem \(\Omega\) pro \(t=K+\kappa\) obtinere potest, fieri \(=M+\theta\theta\omega=M+\frac{\kappa\kappa}{\omega}\). Vice versa, si \(\Omega\) debet valorem aliquem praescriptum \(M+\mu\mu\) non superare, valor ipsius \(t\) necessario inter limites \(K-\mu\BareRoot{\omega}\) et \(K+\mu\BareRoot{\omega}\) contentus esse debet, ubi \(\mu\BareRoot{\omega}\) aequalis fit errori medio in determinatione maxime plausibili ipsius \(t\) metuendo, si pro \(\mu\) accipitur error medius observationum, quibus pondus\(=1\) tribuitur.

\Article{31}
Quoties multitudo quantitatum \(x,y,z\) etc. paullo maior est, determinatio numerica valorum \(A,B,C\) etc. ex aequationibus \(\xi=0,\allowbreak\eta=0.\quad\zeta=0\) etc. per eliminationem vulgarem satis molesta evadit. Propterea in Theoria Motus Corporum Coelestium art.\ 182 algorithmum peculiarem addigitavimus, atque in \textit{Disquisitione de elementis ellipticis Palladis} (Comm. recent. Soc. Gotting Vol.\ I) copiose explicavimus, per quem labor ille ad tantam quantam quidem res fert simplicitatem evehitur. Reducenda scilicet est functio \(\Omega\) sub formam talem:
\[
\frac{u^0u^0}{\mathfrak{A}^0}+\frac{u'u'}{\mathfrak{B}'}+\frac{u''u''}{\mathfrak{C}''}+\frac{u'''u'''}{\mathfrak{D}'''}+\text{etc.}+M
\]
ubi divisores \(\mathfrak{A}^0,\mathfrak{B}',\mathfrak{C}'',\mathfrak{D}'''\) etc. sunt quantitates determinatae; \(u^0,u',u'',u'''\) etc. autem functiones lineares ipsarum \(x,y,z\) etc., quarum tamen secunda \(u'\) libera est ab \(x\), tertia \(u''\) libera ab \(x\) et \(y\), quarta \(u'''\) libera ab \(x,y\) et \(z\), et sic porro, ita ut ultima \(u^{(\pi-1)}\) solam ultimam indeterminatam \(x,y,z\) etc. implicet; denique coëfficientes, per quos \(x,y,z\) etc. resp. multiplicatae sunt in \(u^0,u',u''\) etc., resp. aequales sunt ipsis \(\mathfrak{A}^0,\mathfrak{B}',\mathfrak{C}''\) etc. Quibus ita factis statuendum est \(u^0=0,\allowbreak u'=0,\allowbreak u''=0,\allowbreak u'''=0\) etc., unde valores incognitarum \(x,y,z\) etc. inverso ordine commodissime elicientur. Haud opus videtur, algorithmum ipsum, per quem haec transformatio functionis \(\Omega\) absolvitur, hic denuo repetere.

Sed multo adhuc magis prolixum calculum requirit eliminatio indefinita, cuius adiumento illarum determinationum pondera invenire oportet. Pondus qui\SourceLeaf{043}%
dem determinationis incognitae ultimae (quae sola ultimam \(u^{(\pi-1)}\) ingreditur) per ea, quae in Theoria Motus Corporum Coelestium demonstrata sunt, facile invenitur aequale termino ultimo in serie divisorum \(\mathfrak{A}^0,\mathfrak{B}',\mathfrak{C}''\) etc.; quapropter plures calculatores, ut eliminationem illam molestam evitarent, deficientibus aliis subsidiis, ita sibi consuluerunt, ut algorithmum, de quo diximus pluries, mutato quantitatum \(x,y,z\) etc., ordine, repeterent, singulis deinceps ultimum locum occupantibus. Gratum itaque geometris fore speramus, si modum novum pondera determinationum calculandi, e penitiori argumenti perscrutatione haustum hic exponamus, qui nihil amplius desiderandum relinquere videtur.

\Article{32}
Statuamus itaque esse (I)
\[
\begin{array}{rclllll}
u^0&=&\mathfrak{A}^0x&+\mathfrak{B}^0y&+\mathfrak{C}^0z&+\text{etc.}&+\mathfrak{L}^0\\
u'&=&&\mathfrak{B}'y&+\mathfrak{C}'z&+\text{etc.}&+\mathfrak{L}'\\
u''&=&&&\mathfrak{C}''z&+\text{etc.}&+\mathfrak{L}''\\
\text{etc.}
\end{array}
\]
Hinc erit indefinite
\[
\begin{aligned}
\tfrac12d\Omega&=\xi\,dx+\eta\,dy+\zeta\,dz+\text{etc.}\\[3pt]
&=\frac{u^0du^0}{\mathfrak{A}^0}+\frac{u'du'}{\mathfrak{B}'}+\frac{u''du''}{\mathfrak{C}''}+\text{etc.}\\[3pt]
&=u^0\left(dx+\frac{\mathfrak{B}^0}{\mathfrak{A}^0}dy+\frac{\mathfrak{C}^0}{\mathfrak{A}^0}dz+\text{etc.}\right)\\
&\qquad+u'\left(dy+\frac{\mathfrak{C}'}{\mathfrak{B}'}dz+\text{etc.}\right)+u''(dz+\text{etc.})+\text{etc.}
\end{aligned}
\]
unde colligimus (II)
\[
\begin{aligned}
\xi&=u^0\\
\eta&=\frac{\mathfrak{B}^0}{\mathfrak{A}^0}u^0+u'\\
\zeta&=\frac{\mathfrak{C}^0}{\mathfrak{A}^0}u^0+\frac{\mathfrak{C}'}{\mathfrak{B}'}u'+u''\\
&\text{etc.}
\end{aligned}
\]
Supponamus, hinc derivari formulas sequentes (III)
\SourceLeaf{044}
\[
\begin{aligned}
u^0&=\xi\\
u'&=A'\xi+\eta\\
u''&=A''\xi+B''\eta+\zeta\\
&\text{etc.}
\end{aligned}
\]
Iam e differentiali completo aequationis
\[
\Omega=\xi(x-A)+\eta(y-B)+\zeta(z-C)+\text{etc.}+M
\]
subtracta aequatione
\[
\tfrac12d\Omega=\xi\,dx+\eta\,dy+\zeta\,dz+\text{etc.}
\]
sequitur
\[
\tfrac12d\Omega=(x-A)d\xi+(y-B)d\eta+(z-C)d\zeta+\text{etc.}
\]
quae expressio identica esse debet cum hac ex III demanante:
\[
\begin{aligned}
&\frac{u^0}{\mathfrak{A}^0}.d\xi+\frac{u'}{\mathfrak{B}'}(A'd\xi+d\eta)\\
&\qquad+\frac{u''}{\mathfrak{C}''}(A''d\xi+B''d\eta+d\zeta)+\text{etc.}
\end{aligned}
\]
Hinc colligimus (IV)
\[
\begin{array}{rcllll}
x&=&\dfrac{u^0}{\mathfrak{A}^0}&+A'.\dfrac{u'}{\mathfrak{B}'}&+A''.\dfrac{u''}{\mathfrak{C}''}&+\text{etc.}+A\\[7pt]
y&=&&\dfrac{u'}{\mathfrak{B}'}&+B''.\dfrac{u''}{\mathfrak{C}''}&+\text{etc.}+B\\[7pt]
z&=&&&\dfrac{u''}{\mathfrak{C}''}&+\text{etc.}+C\\[4pt]
\text{etc.}
\end{array}
\]
Substituendo in his expressionibus pro \(u^0,u',u''\) etc. valores earum ex III depromtos eliminatio indefinita absoluta erit. Et quidem ad pondera determinanda habebimus (V)
\[
\begin{aligned}
\relax[\alpha\alpha]&=\frac1{\mathfrak{A}^0}+\frac{A'A'}{\mathfrak{B}'}+\frac{A''A''}{\mathfrak{C}''}+\frac{A'''A'''}{\mathfrak{D}'''}+\text{etc.}\\[4pt]
[\beta\beta]&=\phantom{\frac1{\mathfrak{A}^0}+{}}\frac1{\mathfrak{B}'}+\frac{B''B''}{\mathfrak{C}''}+\frac{B'''B'''}{\mathfrak{D}'''}+\text{etc.}\\[4pt]
[\gamma\gamma]&=\phantom{\frac1{\mathfrak{A}^0}+\frac1{\mathfrak{B}'}+{}}\frac1{\mathfrak{C}''}+\frac{C'''C'''}{\mathfrak{D}'''}+\text{etc.}\\
&\text{etc.}
\end{aligned}
\]
quarum formularum simplicitas nihil desiderandum relinquit. Ceterum etiam pro
\SourceLeaf{045}%
coëfficientibus reliquis \([\alpha\beta],[\alpha\gamma],[\beta\gamma]\) etc. formulae aeque simplices prodeunt, quas tamen, quum illorum usus sit rarior, hic apponere supersedemus.

\Article{33}
Propter rei gravitatem, et ut omnia ad calculum parata sint, etiam formulas explicitas ad determinationem coëfficientium \(A',A'',A'''\) etc., \(B'',B'''\) etc. etc. hic adscribere visum est. Duplici modo hic calculus adornari potest, quum aequationes identicae prodire debeant, tum si valores ipsarum \(u^0,u',u''\) etc. ex III depromti in II substituuntur, tum ex substitutione valorum ipsarum \(\xi,\eta,\zeta\) etc. ex II in III. Prior modus haec formularum systemata subministrat:
\[
\begin{aligned}
\frac{\mathfrak{B}^0}{\mathfrak{A}^0}+A'&=0\\[4pt]
\frac{\mathfrak{C}^0}{\mathfrak{A}^0}+\frac{\mathfrak{C}'}{\mathfrak{B}'}.A'+A''&=0\\[4pt]
\frac{\mathfrak{D}^0}{\mathfrak{A}^0}+\frac{\mathfrak{D}'}{\mathfrak{B}'}.A'+\frac{\mathfrak{D}''}{\mathfrak{C}''}.A''+A'''&=0
\end{aligned}
\]
etc. unde inveniuntur \(A',A'',A'''\) etc.
\[
\begin{aligned}
\frac{\mathfrak{C}'}{\mathfrak{B}'}+B''&=0\\[4pt]
\frac{\mathfrak{D}'}{\mathfrak{B}'}+\frac{\mathfrak{D}''}{\mathfrak{C}''}.B''+B'''&=0
\end{aligned}
\]
etc. unde inveniuntur \(B'',B'''\) etc.
\[
\frac{\mathfrak{D}''}{\mathfrak{C}''}+C'''=0
\]
etc. unde inveniuntur \(C'''\) etc. Et sic porro.

Alter modus has formulas suggerit:
\[
\mathfrak{A}^0A'+\mathfrak{B}^0=0
\]
unde habetur \(A'\).
\[
\begin{array}{rllll}
\mathfrak{A}^0A''&+\mathfrak{B}^0B''&+\mathfrak{C}^0&=0\\
&\mathfrak{B}'B''&+\mathfrak{C}'&=0
\end{array}
\]
unde inveniuntur \(B''\) et \(A''\).
